\section{Related Work}

\paragraph{Appearance and learned visual cues.}
Domain randomization varies simulated appearance to improve transfer
\citep{Tobin2017Domain,Tremblay2018Training}. Texture--shape cue conflicts reveal
that models can rely on visual evidence different from that used by humans
\citep{Geirhos2019Texture}. Our intervention targets an operational source of
appearance variation: changing a scene's material representation while retaining
its geometry. This lets us examine localization and navigation decisions under
a concrete transformation encountered when simulation assets are reused.

\paragraph{Grounding and model readout.}
Shikra and Ferret integrate referring and grounding with language generation
\citep{Chen2023Shikra,You2024Ferret}. Molmo explicitly develops pointing data,
Qwen2.5-VL supports object localization and dynamic image resolution, and
InternVL3 explores multimodal training and inference recipes
\citep{Deitke2024Molmo,Bai2025Qwen,Zhu2025InternVL}. These capabilities make
spatial outputs useful, but the interface through which a point is emitted
also becomes part of the experiment. Prompt dependence is well established for
language models \citep{Zhao2021Calibrate}; we examine its interaction with a
paired visual intervention and separate executable output from conditional
localization accuracy.

\paragraph{Embodied evaluation.}
Vision-and-language navigation couples language understanding to sequential
perception and action \citep{Anderson2018Vision}. GRUtopia provides interactive
scenes and tasks for embodied research \citep{Wang2024GRUtopia}; DualVLN combines
grounding and control in a navigation system \citep{Wei2026Ground}. We use
GRScenes for point-level experiments and KuJiaLe scenes for paired navigation.
The common experimental question is how a change to rendered appearance alters
decisions, from one image to an entire episode.
